\chapter{Functions: The Language of Change}

Calculus is the mathematics of change. But before we can study \emph{how} things
change, we need a precise way to say \emph{what} is changing, and \emph{what it
depends on}. That language is the language of functions. Every derivative is the
rate of change of a function; every integral is the accumulation of a function;
every differential equation is a sentence whose unknown is a function. So this
chapter is not a warm-up to be skimmed. It is the vocabulary of everything that
follows. If you understand functions deeply --- not as formulas to plug into, but
as \emph{rules that turn one quantity into another} --- the rest of calculus will
feel like a story you can already half-tell.

We will build the idea from the ground up: first the intuition (a function is a
machine, a dependable rule), then the precise definition (domain, codomain,
range), then the four ways we meet functions in the wild (formula, table, graph,
words). From there we assemble the toolkit calculus leans on constantly:
composition, inverses, symmetry, monotonicity, the families of polynomial,
rational, and algebraic functions, and the art of transforming a graph by
shifting, stretching, and reflecting. We end by asking the question the whole
chapter is secretly about: \emph{why are functions the objects calculus acts on?}

\section{What a Function Really Is}
\label{ch02:what-is-a-function}

Start with the everyday experience that functions are trying to capture: one
thing depends on another. The cost of filling your tank depends on how many
litres you pump. The distance a dropped stone has fallen depends on how long it
has been falling. The temperature of your coffee depends on how long it has been
sitting. In each case there is an \emph{input} you can dial (litres, seconds) and
an \emph{output} that responds (cost, distance, temperature). A function is the
exact bookkeeper of that dependence.

\begin{intuitionbox}
Picture a machine with a funnel on top and a spout on the bottom. You drop a
number in the top; a single number comes out the bottom. The machine is
\emph{deterministic}: drop in the same number twice and you get the same answer
twice. That is the whole idea. A function is a reliable rule that assigns, to each
allowed input, \emph{exactly one} output.
\end{intuitionbox}

The phrase ``exactly one'' is doing quiet but essential work. If a rule could send
the input $3$ sometimes to $9$ and sometimes to $10$, it would be useless for
prediction --- and prediction is the entire point. So we demand single-valuedness:
one input, one output, no exceptions.

\begin{definition}[Function]
\label{ch02:def-function}
A \emph{function} $f$ from a set $D$ to a set $Y$ is a rule that assigns to each
element $x \in D$ exactly one element $f(x) \in Y$. We write
\[
  f : D \To Y, \qquad x \mapsto f(x).
\]
The set $D$ is the \emph{domain} (the allowed inputs). The set $Y$ is the
\emph{codomain} (a set guaranteed to contain all the outputs). The actual set of
outputs,
\[
  \operatorname{range}(f) = \set{ f(x) : x \in D } \subseteq Y,
\]
is the \emph{range} (or image).
\end{definition}

Read the notation out loud so it stops being decoration. ``$f$ colon $D$ arrow
$Y$'' says \emph{$f$ takes things from $D$ and produces things in $Y$}. The
second arrow, the little ``$\mapsto$'' with a bar, is read ``maps to'' and tells
you the actual rule: $x$ maps to $f(x)$. The first arrow is about \emph{sets}; the
second is about a single \emph{element}.

\begin{focusbox}
Three sets travel with every function, and confusing them causes more grief than
any other early mistake.
\begin{itemize}
  \item \textbf{Domain} $D$: where inputs are allowed to live.
  \item \textbf{Codomain} $Y$: a roomy ``target'' set we promise the outputs fall
        inside. It is part of how the function is \emph{declared}.
  \item \textbf{Range}: the outputs that actually occur. Always a subset of the
        codomain, often a strict subset.
\end{itemize}
For $f:\R\To\R$, $f(x)=x^2$, the codomain is all of $\R$, but the range is only
$[0,\infty)$: no real input ever produces a negative output.
\end{focusbox}

\subsection{The domain is part of the function}

Here is a subtlety beginners routinely miss: \emph{a function is its rule
together with its domain}. The two formulas $g(x)=x$ on $\R$ and $h(x)=x^2/x$ are
not the same function, even though $x^2/x$ simplifies to $x$ wherever it is
defined. The second one is forbidden from eating $x=0$, because division by zero
is undefined, so its domain is $\R\setminus\{0\}$. Same rule almost everywhere,
different domains, different functions. Calculus cares about this fiercely: the
behavior of a function \emph{near} a missing point is exactly the kind of thing a
limit is built to examine.

When a function is handed to you as a bare formula with no domain stated, the
convention is to use the \emph{natural domain}: every real number for which the
formula makes sense (no division by zero, no even roots of negatives, no logs of
non-positive numbers, and so on).

\begin{example}[Finding the natural domain]
\label{ch02:ex-natural-domain}
Find the natural domain of
\[
  f(x) = \frac{\sqrt{x+2}}{\,x-3\,}.
\]
Two requirements must both hold. First, the square root needs a non-negative
argument: $x+2 \ge 0$, i.e. $x \ge -2$. Second, the denominator must not vanish:
$x - 3 \ne 0$, i.e. $x \ne 3$. Intersecting the two conditions gives the natural
domain
\[
  [-2, 3) \cup (3, \infty).
\]
Notice how the domain is a \emph{geometric} object on the number line --- two
intervals with a single point punched out --- not just an algebraic afterthought.
\end{example}

\begin{trapbox}
``$f(x) = x$'' and ``$f(x) = \dfrac{x^2}{x}$'' are different functions because they
have different domains. More generally, cancelling a factor can silently enlarge
the domain. Algebra that ``simplifies'' $\frac{(x-1)(x+2)}{x-1}$ to $x+2$ has
quietly invited $x=1$ back in --- but the original expression was never defined
there. Keep track of what you threw away; in calculus those punched-out points are
where the interesting things (removable discontinuities, limits) live.
\end{trapbox}

\begin{happeningbox}
What is a function ``really,'' underneath the machine metaphor? Mathematicians who
wanted no metaphors at all settled on this: a function \emph{is} its complete
input-output ledger. Formally, $f$ is the set of ordered pairs
\[
  \set{ (x, f(x)) : x \in D },
\]
subject to one rule --- no two pairs share the same first coordinate (that is
single-valuedness restated). So a function is literally a set of points: precisely
the set of points you plot to draw its graph. This is why ``the function'' and
``the graph of the function'' are, at the deepest level, the same object. The
graph is not a picture \emph{of} the function; it \emph{is} the function, drawn.
\end{happeningbox}

\begin{historybox}
The word ``function'' enters mathematics with Leibniz in the 1670s, originally
meaning a quantity (like a tangent length or a curvature) that varied along a
curve. For the next century a ``function'' meant, in practice, a single analytic
formula --- Euler in 1748 defined it as ``an analytic expression composed of a
variable and constants.'' That formula-bound view cracked in the early 1800s when
Fourier showed that wildly different formulas could describe the same curve, and
that strange, piece-by-piece rules deserved to be functions too. Dirichlet in 1837
gave the modern liberation: $y$ is a function of $x$ if to each $x$ there
corresponds a definite $y$ --- \emph{no formula required}. He even offered the
deliberately formula-less example that is $1$ on rationals and $0$ on irrationals.
That shift, from ``function = formula'' to ``function = any single-valued rule,''
is one of the great quiet revolutions in mathematics, and it is exactly the
generality calculus needs.
\end{historybox}

\section{Graphs: Seeing a Function}
\label{ch02:graphs}

If a function is a set of input-output pairs, we can \emph{see} it by plotting each
pair $(x, f(x))$ as a point in the plane, with the input along the horizontal axis
and the output up the vertical axis. Sweep $x$ across the domain and the points
trace a curve: the graph. The graph turns an abstract rule into geometry, and
geometry is where intuition lives.

\begin{figure}[H]
\centering
\begin{tikzpicture}
\begin{axis}[
  width=11cm, height=7cm,
  axis lines=middle,
  xlabel={$x$ (input)}, ylabel={$y=f(x)$ (output)},
  xmin=-3.2, xmax=3.2, ymin=-1.2, ymax=9.2,
  xtick={-3,-2,-1,1,2,3}, ytick={2,4,6,8},
  samples=120, domain=-3:3,
  enlargelimits=false,
  clip=false,
]
\addplot[curveblue, very thick] {x^2};
\addplot[only marks, mark=*, mark size=2pt, tangentred]
  coordinates {(2,4)};
\draw[dashed, gray] (axis cs:2,0) -- (axis cs:2,4);
\draw[dashed, gray] (axis cs:0,4) -- (axis cs:2,4);
\node[anchor=west, tangentred] at (axis cs:2.05,4.6) {$(2,\,f(2))=(2,4)$};
\end{axis}
\end{tikzpicture}
\caption{The graph of $f(x)=x^2$. Each input $x$ on the horizontal axis is sent
straight up (or down) to its single output. The point $(2,4)$ records ``the
machine sends $2$ to $4$.''}
\label{ch02:fig-parabola}
\end{figure}

Figure~\ref{ch02:fig-parabola} also makes the single-valuedness rule visible, and
gives us a quick diagnostic.

\begin{focusbox}
\textbf{The Vertical Line Test.}\quad
A curve in the plane is the graph of a function (of $x$) if and only if every
vertical line crosses it \emph{at most once}. A vertical line is the set of all
points with a fixed input $x$; if it hit the curve twice, that single input would
have two outputs, breaking the ``exactly one'' rule. So the circle $x^2+y^2=1$ is
\emph{not} the graph of a function --- the vertical line $x=0$ stabs it at both
$(0,1)$ and $(0,-1)$.
\end{focusbox}

Learning to read a graph pays off constantly, because calculus keeps asking you
to read off rates and areas from pictures. A graph shows you at a glance where the
function is positive or negative (above or below the $x$-axis), where it is rising
or falling, where it peaks or bottoms out, how fast it climbs (steepness), and what
it does far out toward the edges of the domain. Each of those visual features will
get a precise calculus name in later chapters --- zero, increasing, maximum,
derivative, limit at infinity. The graph is the bridge between seeing and
computing.

\section{Four Ways to Represent a Function}
\label{ch02:four-representations}

The same function can show up wearing four different outfits, and a fluent reader
of calculus can translate among them without thinking. Each representation makes
some things obvious and hides others.

\begin{enumerate}
  \item \textbf{A formula (symbolic).} $f(x) = x^2$. Compact, exact, perfect for
        algebra and for the differentiation rules to come. But a formula can hide
        behavior --- you may not \emph{see} the minimum until you think.
  \item \textbf{A table (numerical).} A list of input-output pairs. This is what a
        thermometer log or an experiment actually produces. Exact about the real
        data; silent about the values \emph{between} the rows.
  \item \textbf{A graph (visual).} The curve. Reveals shape, trend, symmetry, and
        speed at a glance; less good for reading off an exact value.
  \item \textbf{Words (verbal).} ``The output is the square of the input.''
        How problems first arrive in the world. Often the hardest step in an
        applied problem is translating words into one of the other three forms.
\end{enumerate}

\begin{labbox}
Take $f(x) = x^2$ and build all four views by hand. Words: \emph{square the
input}. Formula: $f(x)=x^2$. Table: compute $f$ at $x=-2,-1,0,1,2$ to get
$4,1,0,1,4$. Graph: plot those five points and connect them with a smooth U.
Now ask yourself which representation made the symmetry $f(-x)=f(x)$ jump out,
which made the minimum at $x=0$ obvious, and which would let you read off
$f(1.5)$ most reliably. Different jobs, different outfits.
\end{labbox}

\begin{table}[H]
\centering
\caption{A function as a table: $f(x)=x^2$ sampled at five inputs. The table is
exact at these points and says nothing about what happens between them --- that
silence is precisely what a graph (or a limit) fills in.}
\label{ch02:tab-square}
\begin{tabular}{r r r r r r}
\toprule
$x$    & $-2$ & $-1$ & $0$ & $1$ & $2$ \\
\midrule
$f(x)$ & $4$  & $1$  & $0$ & $1$ & $4$ \\
\bottomrule
\end{tabular}
\end{table}

\begin{applicationbox}
In data science and machine learning, you almost never meet a function as a clean
formula. You meet it as a \emph{table}: rows of measured inputs and outputs (pixels
and labels, features and prices). The central act of supervised learning is to
\emph{guess a formula or graph} that fits the table and generalizes to inputs you
have not seen. The four representations are not academic --- choosing which one to
trust, and how to interpolate between table rows, is the daily work. Calculus
enters the moment you ask ``which way should I nudge the formula to fit the table
better?'' --- that nudge is a derivative.
\end{applicationbox}

\section{Composition: Functions in Series}
\label{ch02:composition}

Machines can be plumbed together: send the output of one into the input of the
next. If $g$ turns $x$ into $g(x)$, and $f$ then turns \emph{that} into
$f(g(x))$, the combined two-stage machine is the \emph{composition} $f \circ g$.

\begin{definition}[Composition]
\label{ch02:def-composition}
Given $g : A \To B$ and $f : B \To C$, their composition $f \circ g : A \To C$ is
defined by
\[
  (f \circ g)(x) = f\bigl(g(x)\bigr).
\]
It is defined exactly at those $x$ in the domain of $g$ for which $g(x)$ lands in
the domain of $f$. You apply $g$ \emph{first}, then $f$ --- read the notation from
the inside out.
\end{definition}

\begin{figure}[H]
\centering
\begin{tikzpicture}[>=Stealth, node distance=2.0cm]
\node[draw, rounded corners, fill=sourceblue!12, minimum width=1.4cm, minimum height=1cm] (x) {$x$};
\node[draw, rounded corners, fill=green!10, minimum width=1.4cm, minimum height=1cm, right=2.4cm of x] (gx) {$g(x)$};
\node[draw, rounded corners, fill=quotientorange!15, minimum width=1.7cm, minimum height=1cm, right=2.6cm of gx] (fgx) {$f(g(x))$};
\draw[->, very thick, sourceblue] (x) -- node[above]{$g$} (gx);
\draw[->, very thick, quotientorange] (gx) -- node[above]{$f$} (fgx);
\draw[->, thick, dashed, structurepurple]
  (x.south) to[bend right=28] node[below]{$f\circ g$} (fgx.south);
\end{tikzpicture}
\caption{Composition is wiring two machines in series. The input $x$ passes through
$g$, and its output is immediately fed to $f$. The dashed arrow is the combined
machine $f \circ g$.}
\label{ch02:fig-composition}
\end{figure}

\begin{example}[Computing a composition]
\label{ch02:ex-composition}
Let $f(u) = \sqrt{u}$ and $g(x) = x^2 + 1$. Then
\[
  (f \circ g)(x) = f(g(x)) = \sqrt{x^2 + 1},
\]
which is defined for \emph{all} real $x$ (since $x^2+1 \ge 1 > 0$ always). Reverse
the order and you get a genuinely different function:
\[
  (g \circ f)(x) = g(f(x)) = (\sqrt{x})^2 + 1 = x + 1,
\]
but now the domain is only $x \ge 0$, inherited from the inner $\sqrt{x}$. So even
where $g\circ f$ and the bare formula $x+1$ agree, the domain remembers the square
root. \textbf{Order matters, and so does the inner domain.}
\end{example}

\begin{trapbox}
$f \circ g$ is almost never the same as $g \circ f$. Composition is not
commutative. Above, one branch gave $\sqrt{x^2+1}$ and the other gave $x+1$ ---
not even close. Always apply the inner function first, and carry its domain
restrictions along.
\end{trapbox}

\begin{aestheticbox}
Why fuss over composition this early? Because the single most important rule in
all of differential calculus --- the \emph{chain rule} --- is the rule for
differentiating a composition. ``Most functions you will ever differentiate are
compositions in disguise'' is barely an exaggeration: $\sin(x^2)$, $e^{3x}$,
$\sqrt{1-x^2}$, $\ln(\cos x)$ are all $f(g(x))$ for some inner $g$. Learning to
\emph{see} the inner and outer machine now is learning to use the chain rule later.
We will return to this exact picture when the chain rule arrives.
\end{aestheticbox}

\section{Inverses: Running the Machine Backwards}
\label{ch02:inverses}

Sometimes you know the output and want the input. You measured the cost; how many
litres did you pump? You know the area of a square; what was the side length? You
want to \emph{undo} the function. The undoing, when it exists, is the
\emph{inverse function} $f^{-1}$.

\begin{intuitionbox}
If $f$ is a machine that sends $x \mapsto y$, the inverse $f^{-1}$ is the machine
that sends $y \mapsto x$ --- it runs the original backwards. Feeding a number
through $f$ and then through $f^{-1}$ must return you to where you started:
$f^{-1}(f(x)) = x$. The inverse exactly cancels $f$.
\end{intuitionbox}

But not every machine can be run backwards. Suppose $f(2) = 4$ and $f(-2) = 4$
(as for $f(x)=x^2$). Hand the inverse the output $4$ and ask for the input --- it
cannot decide between $2$ and $-2$. The inverse would have to be two-valued, which
disqualifies it from being a function at all. The fix is to demand that $f$ never
sends two different inputs to the same output.

\begin{definition}[One-to-one and inverse]
\label{ch02:def-inverse}
A function $f : D \To Y$ is \emph{one-to-one} (injective) if distinct inputs give
distinct outputs:
\[
  x_1 \ne x_2 \implies f(x_1) \ne f(x_2).
\]
If $f$ is one-to-one, its \emph{inverse} $f^{-1}$ has domain $\operatorname{range}(f)$ and
is defined by the rule
\[
  f^{-1}(y) = x \quad\text{exactly when}\quad f(x) = y.
\]
It satisfies $f^{-1}(f(x)) = x$ for all $x\in D$ and $f(f^{-1}(y)) = y$ for all
$y$ in the range.
\end{definition}

\begin{trapbox}
The notation $f^{-1}$ means \emph{inverse function}, not reciprocal. So
$f^{-1}(x) \ne \dfrac{1}{f(x)}$. The reciprocal of $\sin x$ is $\csc x$; the
inverse of $\sin x$ (on a suitable domain) is $\arcsin x$. Same superscript,
completely different objects. Blame history for the clash, and read carefully.
\end{trapbox}

There is a beautiful diagnostic for ``one-to-one,'' the mirror image of the
vertical line test.

\begin{focusbox}
\textbf{The Horizontal Line Test.}\quad
$f$ is one-to-one --- and so has an inverse --- exactly when every \emph{horizontal}
line crosses its graph at most once. A horizontal line is a fixed output value $y$;
if it hit the graph twice, two inputs would share that output, killing
invertibility. $f(x)=x^2$ fails this test on all of $\R$ (the line $y=4$ hits at
$x=\pm 2$), which is exactly why $\sqrt{\,\cdot\,}$ is only the inverse of $x^2$
\emph{after} we restrict to $x \ge 0$.
\end{focusbox}

\subsection{The graph of an inverse is a reflection}

Here is the geometric punchline, and it is genuinely lovely. To invert a function
you swap the roles of input and output --- you swap $x$ and $y$. On a graph,
swapping the two coordinates is exactly reflecting across the diagonal line $y=x$.
So the graph of $f^{-1}$ is the mirror image of the graph of $f$ in that diagonal.

\begin{figure}[H]
\centering
\begin{tikzpicture}
\begin{axis}[
  width=8.6cm, height=8.6cm,
  axis lines=middle,
  xlabel={$x$}, ylabel={$y$},
  xmin=-0.4, xmax=4.2, ymin=-0.4, ymax=4.2,
  xtick={1,2,3,4}, ytick={1,2,3,4},
  samples=100,
  enlargelimits=false,
  clip=false,
]
\addplot[curveblue, very thick, domain=0:2] {x^2};
\addplot[targetgreen, very thick, domain=0:4] {sqrt(x)};
\addplot[gray, dashed, domain=0:4] {x};
\node[curveblue, anchor=west] at (axis cs:1.45,3.3) {$f(x)=x^2$};
\node[targetgreen, anchor=west] at (axis cs:3.1,1.45) {$f^{-1}(x)=\sqrt{x}$};
\node[gray, anchor=west, rotate=45] at (axis cs:2.7,2.9) {$y=x$};
\addplot[only marks, mark=*, mark size=1.6pt, black] coordinates {(1.5,2.25) (2.25,1.5)};
\draw[gray, dotted] (axis cs:1.5,2.25) -- (axis cs:2.25,1.5);
\end{axis}
\end{tikzpicture}
\caption{An inverse is a reflection across $y=x$. The point $(1.5,2.25)$ on
$f(x)=x^2$ (restricted to $x\ge 0$) is mirrored to $(2.25,1.5)$ on
$f^{-1}(x)=\sqrt{x}$. Swapping input and output \emph{is} reflecting in the
diagonal.}
\label{ch02:fig-inverse}
\end{figure}

\begin{whybox}
Why does inverting a function mean reflecting its graph across $y=x$? Trace it
slowly. A point on the graph of $f$ is a pair $(a, b)$ with $b = f(a)$ --- ``input
$a$, output $b$.'' The inverse reverses the roles: it takes input $b$ to output
$a$, so the corresponding point on the graph of $f^{-1}$ is $(b, a)$. The
operation $(a,b) \mapsto (b,a)$ that turns one graph into the other is precisely
the swap of horizontal and vertical coordinates. And swapping coordinates is a
reflection across the line where the two coordinates are equal --- the line $y=x$.
So the geometry is not a coincidence or a trick to memorize; it is the visible
shadow of the algebra ``swap input and output.''
\end{whybox}

\begin{example}[Finding an inverse algebraically]
\label{ch02:ex-find-inverse}
Find the inverse of $f(x) = 2x + 3$. Set $y = 2x + 3$, then \emph{solve for $x$}
(this is the ``swap input and output'' step made concrete):
\[
  y = 2x + 3 \;\implies\; x = \frac{y - 3}{2}.
\]
So $f^{-1}(y) = \dfrac{y-3}{2}$, or, renaming the input back to $x$,
$f^{-1}(x) = \dfrac{x-3}{2}$. Sanity check by composing:
$f^{-1}(f(x)) = \frac{(2x+3)-3}{2} = \frac{2x}{2} = x$. The machine and its
reverse cancel, exactly as required.
\end{example}

\section{Symmetry and Monotonicity: Reading a Function's Character}
\label{ch02:symmetry-monotonicity}

Two qualitative features of a function come up so often that they earn permanent
names. Both are visible at a glance on the graph, and both will matter enormously
once we are differentiating and integrating.

\subsection{Even and odd: the two symmetries}

\begin{definition}[Even and odd functions]
\label{ch02:def-even-odd}
A function $f$ is \emph{even} if $f(-x) = f(x)$ for every $x$ in its domain, and
\emph{odd} if $f(-x) = -f(x)$ for every $x$ in its domain.
\end{definition}

Even means \emph{mirror symmetry across the $y$-axis}: the left half of the graph
is the reflection of the right half. Think $f(x)=x^2$ or $\cos x$. Odd means
\emph{point symmetry through the origin}: rotate the graph $180^\circ$ about
$(0,0)$ and it lands on itself. Think $f(x)=x^3$ or $\sin x$. The names come from
powers of $x$: $x^{\text{even}}$ is an even function, $x^{\text{odd}}$ is an odd
function, which is the cleanest possible mnemonic.

\begin{figure}[H]
\centering
\begin{tikzpicture}
\begin{axis}[
  width=6.6cm, height=5.6cm,
  axis lines=middle, title={even: $f(x)=x^2$},
  xmin=-2.3, xmax=2.3, ymin=-0.5, ymax=4.6,
  xtick=\empty, ytick=\empty, samples=100, domain=-2:2,
  enlargelimits=false,
]
\addplot[curveblue, very thick] {x^2};
\addplot[only marks, mark=*, mark size=1.5pt, black] coordinates {(1.3,1.69) (-1.3,1.69)};
\draw[gray, dashed] (axis cs:-1.3,1.69) -- (axis cs:1.3,1.69);
\end{axis}
\end{tikzpicture}
\hfill
\begin{tikzpicture}
\begin{axis}[
  width=6.6cm, height=5.6cm,
  axis lines=middle, title={odd: $f(x)=x^3$},
  xmin=-1.7, xmax=1.7, ymin=-4.2, ymax=4.2,
  xtick=\empty, ytick=\empty, samples=100, domain=-1.55:1.55,
  enlargelimits=false,
]
\addplot[curveblue, very thick] {x^3};
\addplot[only marks, mark=*, mark size=1.5pt, black] coordinates {(1.2,1.728) (-1.2,-1.728)};
\draw[gray, dashed] (axis cs:-1.2,-1.728) -- (axis cs:1.2,1.728);
\end{axis}
\end{tikzpicture}
\caption{Even functions are mirror-symmetric across the $y$-axis ($f(-x)=f(x)$);
odd functions are rotationally symmetric through the origin ($f(-x)=-f(x)$).}
\label{ch02:fig-even-odd}
\end{figure}

\begin{example}[Testing symmetry algebraically]
\label{ch02:ex-symmetry}
Classify $f(x) = x^4 - 3x^2$. Replace $x$ by $-x$ and simplify:
\[
  f(-x) = (-x)^4 - 3(-x)^2 = x^4 - 3x^2 = f(x).
\]
It equals the original, so $f$ is \emph{even}. Now try $g(x) = x^3 - x$:
\[
  g(-x) = (-x)^3 - (-x) = -x^3 + x = -(x^3 - x) = -g(x),
\]
so $g$ is \emph{odd}. Most functions are neither: $h(x) = x^2 + x$ gives
$h(-x) = x^2 - x$, which is neither $h(x)$ nor $-h(x)$.
\end{example}

\begin{connectionbox}
Symmetry is not just pretty --- it does real computational work. When you integrate
an \emph{odd} function over a symmetric interval $[-a,a]$, the left half exactly
cancels the right half and the integral is $0$, for free, with no antiderivative
needed. For an \emph{even} function the two halves are equal, so
$\int_{-a}^{a} f = 2\int_0^a f$, halving the work. We will use these shortcuts in
the integration chapters; they are gifts that symmetry hands you. Even and odd are
also the building blocks of Fourier analysis, where every function is split into
its even (cosine) and odd (sine) parts.
\end{connectionbox}

\subsection{Increasing and decreasing}

\begin{definition}[Increasing and decreasing]
\label{ch02:def-monotone}
A function $f$ is \emph{increasing} on an interval $I$ if larger inputs give larger
outputs there:
\[
  x_1 < x_2 \implies f(x_1) < f(x_2) \quad\text{for all } x_1, x_2 \in I.
\]
It is \emph{decreasing} on $I$ if $x_1 < x_2 \implies f(x_1) > f(x_2)$. A function
that is one or the other throughout $I$ is called \emph{monotonic} on $I$.
\end{definition}

On the graph, increasing means the curve goes \emph{uphill} as you read left to
right; decreasing means \emph{downhill}. Most interesting functions do both, in
different stretches: $f(x)=x^2$ falls on $(-\infty,0]$ and rises on $[0,\infty)$,
turning around at the bottom of the U.

\begin{figure}[H]
\centering
\begin{tikzpicture}
\begin{axis}[
  width=11cm, height=5.6cm,
  axis lines=middle, xlabel={$x$}, ylabel={$y$},
  xmin=-2.4, xmax=2.4, ymin=-0.4, ymax=4.4,
  xtick={0}, ytick=\empty, samples=120, domain=-2:2,
  enlargelimits=false,
  clip=false,
]
\addplot[curveblue, very thick, domain=-2:0] {x^2};
\addplot[tangentred, very thick, domain=0:2] {x^2};
\node[curveblue, anchor=south] at (axis cs:-1.35,2.4) {decreasing};
\node[tangentred, anchor=south] at (axis cs:1.35,2.4) {increasing};
\addplot[only marks, mark=*, mark size=1.8pt, black] coordinates {(0,0)};
\node[anchor=north] at (axis cs:0,-0.05) {turning point};
\end{axis}
\end{tikzpicture}
\caption{The parabola $f(x)=x^2$ decreases (downhill) on $(-\infty,0]$ and
increases (uphill) on $[0,\infty)$, turning around at the bottom of the U. The
sign of the slope --- and later, of the derivative --- is exactly what flips at the
turning point.}
\label{ch02:fig-monotonicity}
\end{figure}

\begin{connectionbox}
Notice the quiet link to invertibility from the previous section: a strictly
monotonic function automatically passes the horizontal line test, so \emph{every
strictly increasing or strictly decreasing function has an inverse}. Monotonicity
is the cleanest guarantee of invertibility we have. And in the chapters ahead, the
sign of the \emph{derivative} will be precisely the tool that detects where a
function is increasing or decreasing --- positive slope means uphill. So this
informal idea graduates into a sharp, computable test.
\end{connectionbox}

\section{Families of Functions}
\label{ch02:families}

Calculus is easier when you recognize the species of function in front of you,
because whole families share behavior. We build them up from the simplest.

\subsection{Polynomials}

A \emph{polynomial} is built from $x$ using only addition, subtraction, and
multiplication --- no division by $x$, no roots, no exotic operations. Concretely,
\[
  p(x) = a_n x^n + a_{n-1} x^{n-1} + \cdots + a_1 x + a_0,
\]
where the coefficients $a_i$ are constants and $n$ (the highest power with
$a_n\ne 0$) is the \emph{degree}. The natural domain of every polynomial is all of
$\R$: you can always add and multiply real numbers, so nothing is ever forbidden.
Low degrees have familiar shapes --- degree $0$ is a horizontal line, degree $1$ a
slanted line, degree $2$ a parabola, degree $3$ an S-shaped cubic.

\begin{figure}[H]
\centering
\begin{tikzpicture}
\begin{axis}[
  width=12cm, height=6.4cm,
  axis lines=middle, xlabel={$x$}, ylabel={$y$},
  xmin=-2.6, xmax=2.6, ymin=-4.5, ymax=6.5,
  samples=150, domain=-2.4:2.4,
  legend pos=north west, legend cell align=left,
  enlargelimits=false,
]
\addplot[curveblue, very thick] {x};
\addplot[targetgreen, very thick] {x^2 - 2};
\addplot[quotientorange, very thick] {x^3 - 2*x};
\legend{degree $1$: $x$, degree $2$: $x^2-2$, degree $3$: $x^3-2x$}
\end{axis}
\end{tikzpicture}
\caption{Polynomials of degree $1$, $2$, and $3$. Degree controls the maximum
number of turns: a degree-$n$ polynomial can change direction at most $n-1$ times,
which is why the line never turns, the parabola turns once, and the cubic turns
twice.}
\label{ch02:fig-polynomials}
\end{figure}

\begin{aestheticbox}
Polynomials are the friendliest functions in all of calculus, for three reasons.
They are defined everywhere; they are smooth with no breaks or corners; and you can
differentiate or integrate them with one trivial rule applied term by term. The
deeper payoff comes much later: \emph{Taylor's
theorem} will show that a huge class of functions can be approximated, as closely
as you like, by polynomials. So when we eventually want to understand $\sin x$ or
$e^x$ or $\ln x$ near a point, our strategy will be to replace them, locally, with
a polynomial we already understand. Time spent on polynomials now is repaid then.
\end{aestheticbox}

\subsection{Rational functions}

Allow division by polynomials and you get the \emph{rational functions}: ratios
\[
  r(x) = \frac{p(x)}{q(x)}, \qquad p, q \text{ polynomials}, \quad q \not\equiv 0.
\]
The name is the same ``ratio'' as in rational \emph{numbers} (which are ratios of
integers); here it is a ratio of polynomials. The new and important feature is
that the domain now has holes: $r(x)$ is undefined wherever the denominator
vanishes. Those forbidden points are where the most interesting calculus behavior
lives --- vertical asymptotes, where the function blows up toward $\pm\infty$.

\begin{figure}[H]
\centering
\begin{tikzpicture}
\begin{axis}[
  width=10.6cm, height=6.4cm,
  axis lines=middle, xlabel={$x$}, ylabel={$y$},
  xmin=-3.2, xmax=3.2, ymin=-4.5, ymax=4.5,
  samples=200,
  restrict y to domain=-12:12,
  enlargelimits=false,
]
\addplot[curveblue, very thick, domain=-3:-0.34] {1/x};
\addplot[curveblue, very thick, domain=0.34:3] {1/x};
\draw[dashed, trapred, thick] (axis cs:0,-4.5) -- (axis cs:0,4.5);
\node[trapred, anchor=west] at (axis cs:0.12,3.9) {vertical asymptote $x=0$};
\node[curveblue, anchor=west] at (axis cs:1.3,1.4) {$r(x)=\dfrac{1}{x}$};
\end{axis}
\end{tikzpicture}
\caption{The rational function $r(x)=1/x$. Its natural domain is
$\R\setminus\{0\}$; as $x\to 0$ the value shoots off to $\pm\infty$, producing a
vertical asymptote. Holes in the domain are where rational functions get
dramatic.}
\label{ch02:fig-rational}
\end{figure}

\subsection{Algebraic functions}

Go one step further and allow roots --- square roots, cube roots, fractional
powers --- in addition to the four arithmetic operations. Anything you can build
from $x$ with $+,-,\times,\div$ and roots is an \emph{algebraic function}. Examples:
$\sqrt{x}$, $\sqrt[3]{x^2+1}$, $\dfrac{x}{\sqrt{1-x^2}}$. The presence of even
roots restricts the domain (you cannot take $\sqrt{\text{negative}}$ over the
reals), which is why $\sqrt{x}$ lives only on $[0,\infty)$ and
$\dfrac{x}{\sqrt{1-x^2}}$ only on $(-1,1)$.

\begin{focusbox}
A useful hierarchy, each family containing the previous:
\[
  \text{polynomials} \subset \text{rational} \subset \text{algebraic}.
\]
Every polynomial is a rational function (with denominator $1$); every rational
function is algebraic (no roots needed). Functions that are \emph{not} algebraic at
all --- that can never be captured by finitely many arithmetic operations and
roots --- are called \emph{transcendental}. These are the headline acts of the next
chapters: $e^x$, $\ln x$, $\sin x$, $\cos x$. They ``transcend'' algebra, and
calculus is the right tool to tame them.
\end{focusbox}

\section{Transformations: Reshaping a Graph}
\label{ch02:transformations}

Once you know the graph of one function, you can get the graphs of a whole crowd
of relatives for free --- by shifting, stretching, and reflecting. Understanding
transformations means you never start a graph from scratch; you start from a shape
you know and \emph{move} it. There are four basic moves, and the key to all of them
is a single distinction: changes \emph{outside} the function act on the
\emph{output} ($y$), and changes \emph{inside} act on the \emph{input} ($x$).

\begin{focusbox}
\textbf{The four transformations.}\quad
Start from $y = f(x)$ and let $c > 0$.
\begin{itemize}
  \item \textbf{Vertical shift:} $f(x) + c$ moves the graph \emph{up} by $c$;
        $f(x) - c$ moves it \emph{down}. (Outside; behaves as you expect.)
  \item \textbf{Horizontal shift:} $f(x - c)$ moves the graph \emph{right} by $c$;
        $f(x + c)$ moves it \emph{left}. (Inside; behaves \emph{opposite} to the
        sign --- the famous surprise.)
  \item \textbf{Vertical stretch/reflect:} $a\,f(x)$ scales the output by $\abs{a}$
        (a stretch if $\abs{a}>1$, a squash if $\abs{a}<1$) and flips across the
        $x$-axis if $a < 0$. (Outside.)
  \item \textbf{Horizontal stretch/reflect:} $f(bx)$ scales the input axis by the
        factor $1/\abs{b}$, so $\abs{b}>1$ \emph{compresses} horizontally and
        $\abs{b}<1$ stretches it, with a flip across the $y$-axis if $b < 0$.
        (Inside; again the counterintuitive direction.)
\end{itemize}
\end{focusbox}

\begin{figure}[H]
\centering
\begin{tikzpicture}
\begin{axis}[
  width=12cm, height=6.6cm,
  axis lines=middle, xlabel={$x$}, ylabel={$y$},
  xmin=-3.4, xmax=4.4, ymin=-1.2, ymax=6.5,
  samples=120, domain=-3.2:4.2,
  legend pos=north east, legend cell align=left,
  enlargelimits=false,
]
\addplot[curveblue, very thick] {x^2};
\addplot[targetgreen, very thick, dashed] {x^2 + 2};
\addplot[quotientorange, very thick] {(x-1)^2};
\legend{$f(x)=x^2$, up $2$: $x^2+2$, right $1$: $(x-1)^2$}
\end{axis}
\end{tikzpicture}
\caption{The base parabola $x^2$ (solid blue), shifted \emph{up} by $2$ via
$x^2+2$ (an outside change), and shifted \emph{right} by $1$ via $(x-1)^2$ (an
inside change). Note that $(x-1)^2$ moves right even though we \emph{subtracted}
--- the inside surprise.}
\label{ch02:fig-transform}
\end{figure}

\begin{whybox}
Why does $f(x-c)$ shift the graph to the \emph{right} when we \emph{subtract}? It
feels backwards, so let us pin down the mechanism instead of memorizing it. Call
the new function $h(x) = f(x-c)$ and ask: where does $h$ achieve the value that
$f$ achieved at some point $x_0$? It does so when its \emph{inside} equals $x_0$,
that is when $x - c = x_0$, i.e. at $x = x_0 + c$. So whatever $f$ did at $x_0$,
the new function $h$ does at $x_0 + c$ --- a position $c$ units farther right. The
graph's features all move right by $c$. The minus sign is there precisely because
the input has to grow by $c$ to ``catch up'' to the value the inside used to reach.
Once you see it as ``the input must increase to compensate,'' the direction stops
being a paradox.
\end{whybox}

\begin{example}[Reading a transformed graph]
\label{ch02:ex-transform}
Describe how to build the graph of $g(x) = -2(x+3)^2 + 1$ from the parabola
$y = x^2$, taking the operations from the inside out:
\begin{enumerate}
  \item $(x+3)^2$: inside change, shift \emph{left} $3$ units.
  \item $2(x+3)^2$: outside multiply, stretch vertically by $2$ (taller, thinner).
  \item $-2(x+3)^2$: the minus reflects across the $x$-axis (the U now opens
        downward).
  \item $-2(x+3)^2 + 1$: outside add, shift \emph{up} $1$ unit.
\end{enumerate}
Result: a downward parabola with its peak (vertex) at $(-3, 1)$, twice as steep as
the original. You read off the whole shape without plotting a single point.
\end{example}

\begin{figure}[H]
\centering
\begin{tikzpicture}
\begin{axis}[
  width=11cm, height=6.6cm,
  axis lines=middle, xlabel={$x$}, ylabel={$y$},
  xmin=-6.4, xmax=2.4, ymin=-5.5, ymax=4.5,
  samples=120, domain=-6.2:2.2,
  restrict y to domain=-6:5,
  legend pos=south east, legend cell align=left,
  enlargelimits=false,
]
\addplot[curveblue, very thick, dashed] {x^2};
\addplot[tangentred, very thick] {-2*(x+3)^2 + 1};
\addplot[only marks, mark=*, mark size=1.8pt, tangentred] coordinates {(-3,1)};
\node[anchor=south, tangentred] at (axis cs:-3,1.1) {vertex $(-3,1)$};
\legend{base $y=x^2$, $g(x)=-2(x+3)^2+1$}
\end{axis}
\end{tikzpicture}
\caption{Building $g(x)=-2(x+3)^2+1$ from $y=x^2$: shift left $3$, stretch
vertically by $2$, reflect across the $x$-axis, then shift up $1$. The vertex
lands at $(-3,1)$ and the parabola opens downward.}
\label{ch02:fig-transform-example}
\end{figure}

\begin{trapbox}
Order can matter when you stack transformations, because some of them do not
commute. ``Stretch then shift'' generally differs from ``shift then stretch.'' A
safe rule: handle everything \emph{inside} the function (affecting $x$) as one
group and everything \emph{outside} (affecting $y$) as another, and within each
group respect the algebraic order of operations. When in doubt, test one point: a
single tracked point will tell you immediately whether you applied the moves
correctly.
\end{trapbox}

\section{Why Functions Are the Objects Calculus Acts On}
\label{ch02:why-functions}

We close by answering the question lurking under the whole chapter. Calculus has
two central verbs --- \emph{differentiate} and \emph{integrate} --- and both take a
function as input and return a function as output. Differentiation asks ``how fast
is this output changing as the input moves?'' and hands you back the
\emph{rate-of-change function}, the derivative. Integration asks ``how much has
accumulated?'' and hands you back the \emph{accumulation function}, the integral.
Neither question even makes sense for a lone number; both are questions about a
\emph{relationship} between a varying input and a responding output. That
relationship is exactly what a function is.

\begin{happeningbox}
Here is the deep reason functions are unavoidable. To talk about \emph{change}, you
need at least two quantities and a dependence between them: you cannot ask ``how
fast'' without ``how fast \emph{relative to what}.'' The most economical object
that packages ``output depends on input'' is a function. So the moment you decide
to study change quantitatively, the function is not a choice --- it is forced on
you. Calculus did not pick functions for convenience; functions are the smallest
accurate description of a varying world, and calculus is what you can \emph{do}
with them. Everything ahead --- limits, derivatives, integrals, series,
differential equations --- is a verb whose noun is a function.
\end{happeningbox}

\begin{walkbox}
On your next walk, pick something changing around you --- your shadow's length as
the afternoon wears on, your speed as you go up and down a hill, the loudness of
traffic as a car passes. For each one, name the input and the output explicitly.
Then ask the two calculus questions in plain words: \emph{how fast is the output
changing right now?} (that is a derivative) and \emph{how much has accumulated so
far?} (that is an integral). You will find that once you have named the function,
the calculus questions almost ask themselves. That is the whole game, and you can
already play it.
\end{walkbox}

\begin{historybox}
It is no accident that the rise of the function concept and the rise of calculus
are braided together. Newton thought in terms of \emph{fluents} --- quantities
flowing in time --- and \emph{fluxions}, their rates of flow; his ``functions''
were quantities changing with time, the very picture above. Leibniz, who gave us
the word ``function'' and the notation $\dd y/\dd x$ and $\int$, thought in terms
of relations between varying quantities along a curve. For both founders, the
object under study was already a function in all but the modern name. The careful
set-theoretic definition we gave at the start of this chapter came two centuries
later, once mathematicians like Cauchy, Dirichlet, and Weierstrass insisted on
making the foundations airtight. The intuition came first and worked
brilliantly; the rigor came later and explained \emph{why} it worked. That order
--- intuition leading, rigor following to secure the ground --- is the rhythm of
this entire book.
\end{historybox}

\section{Exercises}
\label{ch02:exercises}

\begin{exercise}[Natural domains]
\label{ch02:exc-domains}
Find the natural domain of each function and express it in interval notation.
\begin{enumerate}
  \item $f(x) = \dfrac{1}{x^2 - 4}$
  \item $g(x) = \sqrt{9 - x^2}$
  \item $h(x) = \dfrac{\sqrt{x+1}}{x - 2}$
  \item $k(x) = \dfrac{1}{\sqrt{x} - 2}$ \quad (watch for \emph{two} forbidden
        situations)
\end{enumerate}
\end{exercise}

\begin{exercise}[Range from the graph]
\label{ch02:exc-range}
For each function, determine the range, and justify it geometrically (what is the
lowest/highest output, and is every value between achieved?).
\begin{enumerate}
  \item $f(x) = x^2 - 4$ on $\R$.
  \item $g(x) = \dfrac{1}{x}$ on $(0, \infty)$.
  \item $h(x) = 3 - \abs{x}$ on $\R$.
\end{enumerate}
\end{exercise}

\begin{exercise}[Composition and order]
\label{ch02:exc-composition}
Let $f(x) = x + 1$ and $g(x) = x^2$. Compute $(f\circ g)(x)$ and $(g\circ f)(x)$,
show they are different functions, and find every $x$ at which the two compositions
happen to be \emph{equal}.
\end{exercise}

\begin{exercise}[Decomposing a composition]
\label{ch02:exc-decompose}
Each of the following is a composition $f(g(x))$. Identify a natural inner function
$g$ and outer function $f$. (This is the skill the chain rule will depend on.)
\begin{enumerate}
  \item $\sqrt{x^2 + 1}$
  \item $(3x - 5)^{10}$
  \item $\dfrac{1}{x^2 + 1}$
\end{enumerate}
\end{exercise}

\begin{exercise}[Finding inverses]
\label{ch02:exc-inverses}
Find the inverse of each function on the indicated domain, state the domain of the
inverse, and verify by composition that it undoes the original.
\begin{enumerate}
  \item $f(x) = 3x - 7$ on $\R$.
  \item $g(x) = x^2 + 1$ on $[0, \infty)$.
  \item $h(x) = \dfrac{1}{x - 2}$ on $(2, \infty)$.
\end{enumerate}
\end{exercise}

\begin{exercise}[Even, odd, or neither]
\label{ch02:exc-even-odd}
Classify each function as even, odd, or neither, showing the algebra
$f(-x) = \ldots$ in each case.
\begin{enumerate}
  \item $f(x) = x^4 + 5x^2 - 1$
  \item $g(x) = x^5 - x$
  \item $h(x) = x^3 + x^2$
  \item $p(x) = \dfrac{x}{x^2 + 1}$
\end{enumerate}
\end{exercise}

\begin{exercise}[Transformations]
\label{ch02:exc-transform}
Starting from the base graph $y = x^2$, describe in words (inside-out) the sequence
of transformations that produces each graph, and state the location of the vertex.
\begin{enumerate}
  \item $y = (x - 2)^2 + 3$
  \item $y = -\tfrac{1}{2}(x + 1)^2$
  \item $y = 2(x - 4)^2 - 5$
\end{enumerate}
\end{exercise}

\begin{exercise}[Conceptual: representations]
\label{ch02:exc-representations}
A function is given only by the table
\[
  \begin{array}{c|ccccc}
    x & 0 & 1 & 2 & 3 & 4 \\ \hline
    f(x) & 1 & 2 & 4 & 8 & 16
  \end{array}
\]
Propose a formula that fits all five rows. Explain clearly what the table can and
cannot tell you: in particular, name two \emph{different} functions that agree with
the table at these five inputs but disagree at $x = 1.5$. What does this say about
the limits of the numerical representation?
\end{exercise}

\begin{exercise}[Proof: the vertical line test]
\label{ch02:exc-proof-vlt}
Prove the vertical line test carefully. That is, prove that a set $S$ of points in
the plane is the graph of a function $y = f(x)$ (for $x$ ranging over some domain
$D$) \emph{if and only if} every vertical line meets $S$ in at most one point.
(Hint: a function is, by Definition~\ref{ch02:def-function}, a single-valued
rule; translate ``single-valued'' into a statement about the points $(x,y)\in S$.)
\end{exercise}

\begin{exercise}[Proof: composition of one-to-one functions]
\label{ch02:exc-proof-injective}
Suppose $f$ and $g$ are both one-to-one, and the composition $f \circ g$ is
defined. Prove that $f \circ g$ is also one-to-one. (Hint: assume
$(f\circ g)(x_1) = (f\circ g)(x_2)$ and peel off one function at a time, using
injectivity of each.) Then explain in one sentence why this guarantees $f \circ g$
has an inverse, and state a formula for $(f \circ g)^{-1}$ in terms of $f^{-1}$ and
$g^{-1}$.
\end{exercise}

\begin{exercise}[Proof: even/odd decomposition]
\label{ch02:exc-proof-decomp}
Let $f$ be any function whose domain is symmetric about $0$ (so $-x$ is in the
domain whenever $x$ is). Define
\[
  f_{\text{even}}(x) = \frac{f(x) + f(-x)}{2}, \qquad
  f_{\text{odd}}(x) = \frac{f(x) - f(-x)}{2}.
\]
Prove that $f_{\text{even}}$ is even, that $f_{\text{odd}}$ is odd, and that
$f = f_{\text{even}} + f_{\text{odd}}$. (This shows \emph{every} such function
splits uniquely into an even part plus an odd part --- a fact at the heart of
Fourier analysis.) As a check, carry out the decomposition for $f(x) = e^x$ and
identify the two famous functions you obtain.
\end{exercise}
